% !TEX program = lualatex
% !TeX encoding = UTF-8
\documentclass[12pt,a4paper]{article}

% ============================================================
% إعدادات اللغة العربية (LuaLaTeX + polyglossia)
% ============================================================
\usepackage{polyglossia}
\setmainlanguage{arabic}
\setotherlanguage{english}

% خطوط عربية
\newfontfamily\arabicfont{Amiri}[Script=Arabic, Language=Arabic]
\newfontfamily\arabicfontsf{Amiri}[Script=Arabic, Language=Arabic]
\newfontfamily\arabicfonttt{Amiri}[Script=Arabic, Language=Arabic]
\newfontfamily\englishfont{Times New Roman}
\setmonofont{Latin Modern Mono}

% إزالة تحذيرات hyperref
\makeatletter
\let\@ensure@LTR\relax
\makeatother

% ============================================================
% حزم الرياضيات والتنسيق
% ============================================================
\usepackage{amsmath,amssymb,amsthm}
\usepackage{graphicx}
\usepackage{hyperref}
\usepackage{url}
\usepackage{xcolor}
\usepackage{booktabs}
\usepackage{geometry}
\geometry{margin=2.5cm}

% ============================================================
% أوامر YUCT المخصصة
% ============================================================
\newcommand{\Keff}{K_{\mathrm{eff}}}
\newcommand{\kappac}{\kappa_c}
\newcommand{\eps}{\varepsilon}
\newcommand{\betagamma}{\beta\gamma}
\newcommand{\betac}{\beta}
\newcommand{\tauf}{\tau}

% ============================================================
% بيئات النظريات (بالعربية)
% ============================================================
\newtheorem{theorem}{نظرية}
\newtheorem{lemma}{ليما}
\newtheorem{corollary}{نتيجة طبيعية}
\newtheorem{definition}{تعريف}
\newtheorem{prediction}{تنبؤ}

% ============================================================
% بيانات PDF
% ============================================================
\hypersetup{
	pdftitle={الملحق S: التأخر الزمني السلبي للفوتونات في الذرات الباردة كمظهر من مظاهر الديناميكا الكمومية المنسقة: تحليل منقح في إطار نظرية التنسيق الموحدة ياكوشيف (YUCT)},
	pdfauthor={أليكسي ف. ياكوشيف},
	pdfsubject={YUCT, التأخر الزمني السلبي, الفوتونات, الذرات الباردة, التنسيق الكمومي},
	pdfkeywords={YUCT, التأخر الزمني السلبي, الفوتونات, الذرات الباردة, التنسيق الكمومي, التدرج اللوغاريتمي},
	breaklinks=true,
	pdfencoding=auto,
	bookmarksnumbered=true,
	bookmarksopen=true,
	bookmarksopenlevel=1
}

\begin{document}
	
	\begin{titlepage}
		\centering
		\vspace*{1cm}
		{\Huge\bfseries الملحق S: التأخر الزمني السلبي للفوتونات في الذرات الباردة كمظهر من مظاهر الديناميكا الكمومية المنسقة: \\ تحليل منقح في إطار نظرية التنسيق الموحدة ياكوشيف (YUCT)}
		\vspace{1cm}
		
		{\large أليكسي ف. ياكوشيف\textsuperscript{1}}
		\vspace{0.3cm}
		{\normalsize \textsuperscript{1}Yakushev Research, \\ 
			\textbf نظرية YUCT: \url{https://yuct.org/}, \\ 
			\textbf بروتوكول YPSDC: \url{https://ypsdc.com/}\\
			\textbf DOI: \url{https://doi.org/10.5281/zenodo.18444598}}
		\vspace{1cm}
		{\normalsize مايو 2026 \\ \large الإصدار V.3.0. \\(التدرج اللوغاريتمي المصحح وتراجع التنبؤ بنسبة 13\%)}
		
		\vfill
		
		\begin{abstract}
			\noindent
			أظهرت تجارب حديثة رائدة (Angulo وآخرون، 2024) أن الفوتونات التي تعبر سحابة ذرية فائقة البرودة يمكن أن تخرج من الوسط بتأخر زمني سلبي – حيث تظهر قمة الحزمة الموجية في وقت أبكر مما دخلت فيه. هذه الظاهرة غير البديهية، مع أنها لا تنتهك السببية، تكشف عن تأثيرات تداخل كمومي عميقة. في هذا العمل، نعيد النظر في التفسير في إطار نظرية التنسيق الموحدة ياكوشيف (YUCT). بالبناء على قانون التدرج العام \(\eps = \kappac \alpha (\ln\Keff)^{\betac}\) مع \(\betac = 2/3\) و \(\kappac = 1/3\) (المشتق في الملحق X)، نظهر أن العلاقة الصحيحة بين التأخر الجماعي \(\tau_g\) وكفاءة التنسيق \(\Keff\) هي علاقة لوغاريتمية: \(|\tau_g| = \alpha_0 \ln\Keff\). وجود التأخرات الزمنية السلبية مؤكد بالفعل بالتجربة؛ وقد لوحظت قيم تصل إلى \((-1.14 \pm 0.22)\tau_0\). ومع ذلك، وخلافًا للافتراض الخطي السابق، فإن الشكل اللوغاريتمي يعني أن الاعتماد على درجة الحرارة لـ \(\tau_g\) ضعيف للغاية. باستخدام التدهور الحراري القانوني \(\Keff(T) = K_0 (T/T_0)^{-\gamma}\) مع \(\gamma \approx 0.30\) (المشتق من بيانات جيوفيزيائية)، نجد أن التغير النسبي في \(|\tau_g|\) عند مضاعفة درجة الحرارة هو من رتبة \(10^{-8}\)، أي أقل بكثير من حساسية التجارب الحالية. وبالتالي، فإن التأثير المزعوم بنسبة 13\% هو نتيجة مصطنعة لنموذج خطي خاطئ ويجب التراجع عنه. نقدم الإطار النظري المصحح ونناقش اختبارات تجريبية بديلة يمكن أن تكون حساسة للتدرج اللوغاريتمي، مثل تغيير العمق البصري أو مدة النبضة. يبقى التنبؤ الأساسي لـ YUCT – أن التأخرات الجماعية السلبية هي مظهر من مظاهر كفاءة التنسيق – صحيحًا، لكن قابليته للاختبار في أنظمة الذرات الباردة تتطلب إعدادات محسنة.
		\end{abstract}
		
		\vspace{0.5cm}
		\noindent\textbf{الكلمات المفتاحية:} YUCT، كفاءة التنسيق، التأخر الزمني السلبي، تفاعل فوتون-ذرة، التدرج اللوغاريتمي، الاستقلال عن درجة الحرارة، التراجع.
		
	\end{titlepage}
	
	\tableofcontents
	\newpage
	
	\section{مقدمة}
	
	كان انتشار الضوء عبر الأوساط الرنانة حجر زاوية في البصريات الكمومية لعقود. التأخرات الجماعية، وإعادة تشكيل النبضات، والتأثيرات فائقة السرعة مفهومة جيدًا نظريًا وتجريبيًا \cite{Brillouin1960, Milonni2004}. ومع ذلك، تبقى مسألة أساسية: ما هو المعنى الفيزيائي للزمن الذي يقضيه الفوتون كإثارة ذرية؟
	
	تناولت تجربة رائدة أجراها Angulo وآخرون \cite{Angulo2024} هذه المسألة عن طريق قياس تحول الطور عبر كير الذي يحدثه فوتون منقول على حزمة استقصاء منفصلة. وجدوا أن زمن الإثارة الذرية المتكامل \(\tau_T\) لفوتون منقول يمكن أن يكون سالبًا، مطابقًا للتأخر الجماعي \(\tau_g\) حتى عندما يكون \(\tau_g < 0\). على وجه التحديد، أبلغت التجربة عن قيم \(\tau_T/\tau_0 = -1.14 \pm 0.22\) (لنبضة 10 نانوثانية، عمق بصري OD~4) و \(0.54 \pm 0.28\) (لبارامترات مختلفة) \cite{Angulo2024, Nixon2024}. هذه النتيجة الرائعة تتحدى البديهيات البسيطة وتستدعي إطارًا نظريًا أعمق.
	
	تقدم نظرية التنسيق الموحدة ياكوشيف (YUCT) \cite{Yakushev2026a} مثل هذا الإطار. تفترض YUCT أن جميع التفاعلات الفيزيائية، من البصريات الكمومية إلى الجاذبية، هي مظاهر لعملية تنسيق أساسية تُكمَّن بواسطة \textbf{كفاءة التنسيق} \(\Keff\). يربط قانون تدرج عام أي خطأ نسبي \(\eps\) بـ \(\Keff\) (انظر الملحق X، المعادلة~\ref{eq:error-law-corrected}):
	\begin{equation}
		\eps = \kappac \alpha (\ln\Keff)^{\betac}, \qquad \betac = \frac{2}{3},\ \kappac = \frac{1}{3},
		\label{eq:scaling}
	\end{equation}
	حيث \(\alpha\) يعتمد على النظام. تم التحقق من هذا القانون عبر أكثر من 40 مرتبة مقدارية، من أخطاء تضاعف الحمض النووي إلى تذبذبات الخلفية الميكروية الكونية \cite{Yakushev2026c}. في سياق تفاعل الضوء والمادة، يقيس \(\Keff\) التماسك بين الفوتون والمجموعة الذرية.
	
	نظهر هنا أن تجربة التأخر الزمني السلبي هي نتيجة مباشرة لـ YUCT. ومع ذلك، يجب علينا تصحيح محاولة سابقة افترضت علاقة خطية بين \(\tau_g\) و \(\Keff\). العلاقة الصحيحة، المشتقة من البنية المعلوماتية النظرية لـ YUCT، هي لوغاريتمية. هذا يغير بشكل جذري الاعتماد المتوقع على درجة الحرارة، مما يجعله مهملاً مع الدقة التجريبية الحالية. لذلك نتراجع عن التنبؤ بنسبة 13\% المنشور سابقًا ونقدم تحليلاً منقحًا ومتسقًا ذاتيًا.
	
	\section{المفاهيم الأساسية لـ YUCT}
	
	\subsection{كفاءة التنسيق \(\Keff\)}
	
	في YUCT، يتميز أي نظام متفاعل بقاموس \(D\) (بروتوكولات وحالات مشتركة) ومؤشر \(i\) ينشط إجراءً منسقًا محددًا. كفاءة التنسيق هي
	\begin{equation}
		\Keff = \frac{H(\text{إجراءات})}{H(\text{مؤشر})},
	\end{equation}
	حيث \(H\) تدل على إنتروبيا شانون. بالنسبة للأنظمة الكمومية، يمكن أن يكون \(\Keff\) كبيرًا جدًا، مقتربًا من اللانهاية في حد التشابك المثالي.
	
	\subsection{قانون التدرج العام}
	
	تخضع التقلبات والأخطاء للمعادلة~\eqref{eq:scaling}. ينشأ هذا القانون من الهندسة الكسيرية للقواميس وتم تأكيده تجريبيًا في أنظمة متنوعة مثل تحلل النيوترونات ومنحنيات دوران المجرات \cite{Yakushev2026c}. الشكل اللوغاريتمي ضروري: فهو يعكس حقيقة أن الدقة الفعالة في فضاء المؤشرات تنمو فقط لوغاريتميًا مع حجم القاموس.
	
	\subsection{الاعتماد على درجة الحرارة لـ \(\Keff\)}
	
	في العديد من الأنظمة الفيزيائية، تقل كفاءة التنسيق مع درجة الحرارة بسبب التقلبات الحرارية. من الظواهر الحرجة والتدرج الكسيري، لدينا
	\begin{equation}
		\Keff(T) = K_0 \left(\frac{T}{T_0}\right)^{-\gamma}, \quad \gamma > 0,
		\label{eq:K_T}
	\end{equation}
	حيث \(\gamma\) أس يعتمد على المادة. تم تحديد حاصل الضرب \(\betac\gamma\) بشكل مستقل من بيانات جيوفيزيائية: \(\betagamma = 0.20 \pm 0.03\) \cite{Yakushev2026b}. هذا يعني \(\gamma = \betagamma / \betac = 0.20 / (2/3) = 0.30\). هذه القيمة متسقة مع حجج التدرج العامة للأنظمة غير المنتظمة.
	
	\section{نموذج YUCT لتفاعل الفوتون-الذرة}
	
	لنأخذ نبضة فوتونية رنانة تسقط على مجموعة من الذرات الباردة. تعمل الذرات كـ \textbf{قاموس} \(D\): يتم تحضيرها في حالة كمومية محددة (مثلًا عن طريق التبريد بالليزر والمصيدة). يحمل الفوتون \textbf{مؤشرًا} – شكل حزمته الموجية وتردده. يؤدي التفاعل إلى حالة منسقة حيث قد ينتقل الفوتون أو يمتص أو يشتت.
	
	في تجربة Angulo وآخرون \cite{Angulo2024}، الكمية محل الاهتمام هي \textbf{زمن الإثارة الذرية} \(\tau_T\) لفوتون منقول. من YUCT، هذا الزمن متناسب مع التأخر الجماعي \(\tau_g\). يُظهر تحليل القيمة الضعيفة لـ Thompson وآخرون \cite{Thompson2023} أن \(\tau_T = \tau_g\) لفوتون منقول في النظام ذي الصلة.
	
	الخطوة الحاسمة هي ربط \(\tau_g\) بـ \(\Keff\). ينشأ التأخر الجماعي من التفاعل المتماسك بين الفوتون والمجموعة الذرية. مقدار هذا التفاعل ليس متناسبًا مع \(\Keff\) نفسه (الذي يمكن أن يكون هائلاً) بل مع كمية \textbf{المعلومات} التي يحملها الفوتون عن الحالة الذرية. في YUCT، مقياس المعلومات ذو الصلة هو لوغاريتم كفاءة التنسيق، لأن طول المؤشر \(\ell = \log_2 M\) يحدد عدد الحالات التي يمكن تمييزها. لذلك نفترض:
	\begin{equation}
		|\tau_g| = \alpha_0 \, \ln\Keff,
		\label{eq:tau_log}
	\end{equation}
	حيث \(\alpha_0\) ثابت زمني أساسي له بعد الزمن. لخط الروبيديوم D2، يعطي العرض الطبيعي \(\Gamma = 2\pi \times 6\) MHz زمنًا مميزًا \(\tau_{\text{nat}} = 1/\Gamma \approx 26\) نانوثانية. المطابقة مع القيمة التجريبية \(|\tau_g| \approx 30\) نانوثانية عند \(K_0 \sim 10^4\) تعطي \(\alpha_0 = 30\,\text{ns} / \ln(10^4) = 30 / 9.21 \approx 3.26\) نانوثانية. هذا معقول: زمن التفاعل الأولي هو من رتبة بضع نانوثواني.
	
	\subsection{الاعتماد الحراري للتأخر الجماعي}
	
	بتعويض المعادلة~\eqref{eq:K_T} في المعادلة~\eqref{eq:tau_log} نحصل على
	\begin{equation}
		|\tau_g(T)| = \alpha_0 \ln\!\left(K_0 \left(\frac{T}{T_0}\right)^{-\gamma}\right)
		= \alpha_0 \ln K_0 - \alpha_0 \gamma \ln\!\left(\frac{T}{T_0}\right).
	\end{equation}
	باعتبار \(|\tau_g(T_0)| = \alpha_0 \ln K_0\)، نحصل على
	\begin{equation}
		|\tau_g(T)| = |\tau_g(T_0)| - \alpha_0 \gamma \ln\!\left(\frac{T}{T_0}\right).
		\label{eq:tau_T_log}
	\end{equation}
	هذه هي النتيجة المركزية: مقدار التأخر الجماعي يقل \textbf{لوغاريتميًا} مع زيادة درجة الحرارة. لنسبة درجات حرارة \(T/T_0 = 2\)، يكون التغير
	\begin{equation}
		\Delta |\tau_g| = \alpha_0 \gamma \ln 2.
	\end{equation}
	باستخدام \(\alpha_0 \approx 3.26\) نانوثانية و \(\gamma = 0.30\)، نجد \(\Delta |\tau_g| \approx 3.26 \times 0.30 \times 0.693 \approx 0.68\) نانوثانية. التغير النسبي بالنسبة للقيمة الأساسية \(|\tau_g(T_0)| \approx 30\) نانوثانية هو \(0.68 / 30 \approx 2.3\%\). هذا لا يزال ضمن عدم اليقين الإحصائي للتجربة (\(\pm 0.22\) بوحدات \(\tau_0\)، أي حوالي \(\pm 5.7\) نانوثانية). ومع ذلك، فإن التأثير على حافة قابلية الكشف. الأهم من ذلك، أن التغير هو \textbf{لوغاريتمي}، وليس دالة قوى، ومقداره أصغر بكثير من الـ 13\% التي ادُعيت سابقًا.
	
	\subsection{مقارنة مع النموذج الخطي السابق (الخاطئ)}
	
	في الأعمال السابقة، افترضت علاقة خطية \(|\tau_g| \propto \Keff\). أدى ذلك إلى \(|\tau_g| \propto T^{-\gamma}\)، ومع \(\gamma = 0.30\)، تغير بنسبة 13\% لمضاعفة درجة الحرارة. استند هذا التنبؤ إلى سوء فهم دور \(\Keff\) في الأنظمة الكمومية. تنشأ العلاقة اللوغاريتمية الصحيحة لأن \(\Keff\) هي نسبة إنتروبيات، وليست سعة حقل فيزيائي مباشر. النموذج الخطي غير متوافق مع الأسس المعلوماتية النظرية لـ YUCT ويُتراجع عنه هنا.
	
	\section{فحص الاتساق العددي من البيانات الموجودة}
	
	توفر تجربة Angulo وآخرون \cite{Angulo2024} نقطة بيانات محددة تسمح لنا بتقدير \(\alpha_0\) و \(K_0\). بالنسبة لتكوين النبضة 10 نانوثانية، العمق البصري OD~4، تم الإبلاغ عن زمن إثارة ذرية \(\tau_T = -1.14 \tau_0\). بأخذ \(\tau_0 \approx 26\) نانوثانية، لدينا \(|\tau_g| \approx 30\) نانوثانية. بافتراض أن \(\Keff(T_0) = K_0\) من رتبة \(10^4\) (نموذجي لمجموعة ذرية باردة مع \(10^6\) ذرة وتماسك عالٍ)، نحصل على \(\alpha_0 = 30\,\text{ns} / \ln(10^4) \approx 3.26\) نانوثانية. هذه القيمة معقولة كمقياس زمني مميز للتفاعل الأولي بين الفوتون والذرة. لا حاجة لبارامترات حرة إضافية.
	
	\section{مقارنة مع البصريات الكمومية القياسية}
	
	في الإطار التقليدي للبصريات الكمومية، يُعطى التأخر الجماعي \(\tau_g\) لنبضة رنانة بمشتق طور النفاذية بالنسبة للتردد، \(\tau_g = d\phi/d\omega\). بالنسبة لاستجابة ذرية لورنتزية، يعطي هذا شكلًا مميزًا يمكن أن يصبح سالبًا على جناح الرنين بسبب التشتت الشاذ \cite{Milonni2004}. ومع ذلك، لا تتنبأ النظرية القياسية بأي اعتماد لـ \(\tau_g\) على درجة حرارة السحابة الذرية beyond تأثيرات تافهة مثل الاتساع الدوبلري (الذي هو خطي في درجة الحرارة، وليس لوغاريتميًا). يُفهم التأخر الزمني السلبي المرصود على أنه تأثير تداخل متماسك لا يتضمن تبادل طاقة.
	
	تتجاوز YUCT النظرية القياسية من خلال ربط مقدار التأخر بكفاءة التنسيق \(\Keff\). الاعتماد اللوغاريتمي هو سمة مميزة: إذا كان بإمكان المرء تغيير \(\Keff\) بتغيير الكثافة الذرية، أو العمق البصري، أو مدة النبضة، فيجب أن يتدرج التأخر كـ \(\ln\Keff\). هذا تنبؤ قابل للاختبار لا يعتمد على تغيرات درجة الحرارة.
	
	\section{مقترحات تجريبية منقحة}
	
	لأن الاعتماد على درجة الحرارة ضعيف جدًا (لوغاريتمي، مع تغير نسبي بضع بالمائة فقط لمضاعفة درجة الحرارة)، فمن غير المرجح أن يمكن اكتشافه بالدقة الحالية. لذلك نقترح تجارب بديلة تغير \(\Keff\) بشكل مباشر بطريقة أكثر تحكمًا.
	
	\subsection{تغيير العمق البصري}
	
	كفاءة التنسيق \(\Keff\) متناسبة مع عدد الذرات في المجموعة وقوة التفاعل. بتقليل العمق البصري (مثلًا بخفض الكثافة الذرية)، يقل \(\Keff\). يجب أن يتبع مقدار التأخر الجماعي \(|\tau_g| = \alpha_0 \ln\Keff\). لتخفيض \(\Keff\) بعامل 10 (مثلًا من \(10^4\) إلى \(10^3\))، سيكون التغير في \(|\tau_g|\) هو \(\alpha_0 (\ln 10^4 - \ln 10^3) = \alpha_0 \ln 10 \approx 3.26 \times 2.3026 \approx 7.5\) نانوثانية، وهو حوالي 25\% من القيمة الأساسية ويسهل قياسه. مثل هذه التجربة بسيطة: كرر قياس \(\tau_T\) لكثافات ذرية مختلفة (مع إبقاء جميع البارامترات الأخرى ثابتة) وتحقق من التدرج اللوغاريتمي.
	
	\subsection{تغيير مدة النبضة}
	
	يعتمد المحتوى المعلوماتي لنبضة الفوتون (طول المؤشر) على مدتها الزمنية وعرض نطاقها. تحمل النبضات الأقصر معلومات طيفية أكثر وبالتالي لها \(\Keff\) أكبر. بتغيير مدة النبضة بشكل منهجي من 1 نانوثانية إلى 20 نانوثانية، يمكن ملاحظة كيف يتدرج التأخر الجماعي مع \(\ln(\text{عرض النطاق})\). هذا يوفر اختبارًا مباشرًا آخر للقانون اللوغاريتمي.
	
	\subsection{الاعتماد على درجة الحرارة كاختبار صفري}
	
	على الرغم من أن تأثير درجة الحرارة المتوقع ضئيل، إلا أنه ليس صفرًا تمامًا. قياس صفري (لا تغير يمكن اكتشافه ضمن الأخطاء التجريبية) سيكون متسقًا مع نموذجنا اللوغاريتمي، بينما كان النموذج الخطي السابق قد تنبأ بتغير واضح بنسبة 13\%. وبالتالي، فإن البيانات الموجودة بالفعل لا تؤيد النموذج الخطي، وقد تكتشف تجارب المستقبل عالية الدقة في النهاية التحول اللوغاريتمي. نحن نشجع المجربين على اختبار كليهما.
	
	\section{الارتباط بالتنبؤ بكتلة الفوتون: فحص اتساق}
	\label{sec:photon_mass_consistency}
	
	أنتج إطار YUCT أيضًا تنبؤًا أعمى للكتلة الساكنة للفوتون، مشتقًا من نفس سلم الكتل \(m = m_e q^{N_f}\) مع \(q=(3/2)^{1/3}\). العقدة المثلى المتسقة مع الحدود التجريبية العليا هي \(N_\gamma = -410\)، مما يعطي
	\[
	m_\gamma \approx 7.3 \times 10^{-19}~\text{eV} \quad (\text{انظر الملحق Photon Mass}).
	\]
	
	قد يتساءل المرء ما إذا كانت كتلة فوتون صغيرة جدًا غير صفرية يمكن أن تؤثر على قياس التأخر الجماعي في الذرات الباردة. الطول الموجي لكومبتون الموافق لـ \(m_\gamma\) هو
	\[
	\lambda_c = \frac{\hbar}{m_\gamma c} \approx 2.7\times10^{11}~\text{m},
	\]
	وهو أكبر بكثير من السحابة الذرية (الحجم النموذجي \(\sim 1\) مم). وبالتالي، فإن تأثيرات الانتشار بسبب كتلة فوتون محدودة مهملة تمامًا في هذه التجربة. التأخر الزمني السلبي هو تأثير تنسيق بحت، مستقل عن مقياس الكتلة المطلق.
	
	بالمقابل، فإن حقيقة أن إطارًا واحدًا يتنبأ بشكل صحيح بكل من:
	\begin{itemize}
		\item وجود تأخرات جماعية سلبية بمقدار محدد بواسطة \(\ln\Keff\) (هذا الملحق)،
		\item الحد الأعلى لكتلة الفوتون (\(m_\gamma < 10^{-18}\) eV) وعقدته المنفصلة \(N_\gamma = -410\) (ملحق Photon Mass)،
	\end{itemize}
	تشكل فحص اتساق قوي. يتعلق التنبؤان بمجالات فيزيائية مختلفة تمامًا (البصريات الكمومية مقابل خصائص الجسيمات الأساسية) ومع ذلك يشتركان في نفس الثوابت الأساسية \(q\)، \(S_{\mathrm{odd}}/S_{\mathrm{even}}\)، ومفهوم كفاءة التنسيق \(\Keff\). لا حاجة لتعديلات مخصصة لجعلها متوافقة.
	
	وبالتالي، فإن تجربة التأخر الفوتوني السلبي لا تتعارض مع تنبؤ YUCT لكتلة الفوتون؛ بل يعمل كلاهما كركائز مستقلة تدعم نموذج التنسيق الموحد.
	
	\section{التراجع عن التنبؤ الأعمى السابق}
	
	في إصدار سابق من هذا الملحق (V.2.0)، توقعنا خطأً أن التأخر الجماعي سيتدرج كـ \(T^{-0.20}\)، مما يؤدي إلى تغير بنسبة 13\% عند مضاعفة درجة الحرارة. استند هذا التنبؤ إلى علاقة خطية خاطئة \(|\tau_g| \propto \Keff\). تظهر العلاقة اللوغاريتمية الصحيحة المشتقة هنا أن تأثير درجة الحرارة أصغر بحوالي مرتبتين من حيث الحجم. لذلك، فإن الادعاء بتأثير 13\% غير صحيح ويُتراجع عنه رسميًا. نعتذر عن هذا السهو ونشكر المشاركين في التصحيح الرياضي على تحديد عدم الاتساق.
	
	ومع ذلك، يبقى المفهوم الأساسي لـ YUCT – أن التأخرات الجماعية السلبية هي مظهر من مظاهر كفاءة التنسيق – صحيحًا. يؤكد التأكيد التجريبي للأزمنة الذرية السلبية بالفعل فكرة أن الفوتون والذرة يشتركان في قاموس منسق. يوفر قانون التدرج المصحح الآن إطارًا نظريًا محسنًا للاختبارات المستقبلية.
	
	\section{الاستنتاج}
	
	صححنا تحليل YUCT لتجربة التأخر الفوتوني السلبي. قانون الخطأ العام لـ YUCT هو لوغاريتمي، وبالتالي يجب أن تكون العلاقة بين التأخر الجماعي وكفاءة التنسيق لوغاريتمية: \(|\tau_g| = \alpha_0 \ln\Keff\). هذا يؤدي إلى اعتماد ضعيف جدًا على درجة الحرارة (لوغاريتمي، مع تغير نسبي بضع بالمائة فقط عند مضاعفة درجة الحرارة) ويزيل تأثير 13\% الذي ادُعي سابقًا. يُتراجع عن التنبؤ الخاطئ. تقدم الاختبارات البديلة، مثل تغيير العمق البصري أو مدة النبضة، وسائل أكثر حساسية بكثير للتدرج اللوغاريتمي. يبقى وجود التأخرات الزمنية السلبية تأكيدًا مذهلاً للطبيعة المنسقة للتفاعلات الكمومية، وتوفر YUCT تفسيرًا معلوماتيًا نظريًا متسقًا.
	
	\section*{شكر وتقدير}
	
	يشكر المؤلف المشاركين في التصحيح الرياضي للإشارة إلى عدم الاتساق بين الافتراض الخطي والبنية اللوغاريتمية لـ YUCT. كما يشكر المشاركين في ندوة YUCT على النقاشات المحفزة وفريق تورونتو (Aephraim Steinberg و Daniela Angulo وزملائهم) على تجربتهم الجميلة.
	
	\section*{توفر البيانات}
	
	جميع الحسابات والمواد الإضافية متاحة على \url{https://github.com/Alexey-Yakushev-YUCT/YPSDC}.
	
	\begin{thebibliography}{99}
		
		\bibitem{Angulo2024}
		D. Angulo, K. Thompson, A. Jiao, H. Wiseman, A. Steinberg, ``Experimental evidence that a photon can spend a negative amount of time in an atom cloud,'' arXiv:2409.03680 [quant-ph] (2024).
		
		\bibitem{Nixon2024}
		V.-M. Nixon et al., ``Measuring the time transmitted photons spend as atomic excitations,'' American Physical Society DAMOP Conference, Session Y07 (June 2024).
		
		\bibitem{Yakushev2026a}
		A. V. Yakushev, ``Yakushev Unified Coordination Theory (YUCT),'' Zenodo, \url{https://doi.org/10.5281/zenodo.18444599} (2026).
		
		\bibitem{Yakushev2026b}
		A. V. Yakushev, ``Appendix P: Temperature Dependence of Gravity,'' Zenodo (2026).
		
		\bibitem{Yakushev2026c}
		A. V. Yakushev, ``Appendix L: Fractal Error Scaling,'' Zenodo (2026).
		
		\bibitem{Yakushev2026d}
		A. V. Yakushev, ``Appendix X: Generalization of Shannon Information Theory,'' Zenodo (2026).
		
		\bibitem{Thompson2023}
		K. Thompson et al., ``How much time does a photon spend as an atomic excitation before being transmitted?'' arXiv:2310.00432 (2023).
		
		\bibitem{Brillouin1960}
		L. Brillouin, \textit{Wave Propagation and Group Velocity} (Academic Press, 1960).
		
		\bibitem{Milonni2004}
		P. Milonni, \textit{Fast Light, Slow Light and Left-Handed Light} (CRC Press, 2004).
		
		\bibitem{Wiseman2023}
		H. M. Wiseman, A. M. Steinberg, and M. Hallaji, ``Obtaining a single-photon weak value from experiments using a strong coherent state,'' AVS Quantum Sci. \textbf{5}, 024401 (2023).
		
	\end{thebibliography}
	
\end{document}